%% =====================================================================
%% Self-Harnessing Recursive Language Models — AAAI 2027 version
%% Adapted from proposal.tex (NeurIPS 2026 workshop version).
%% =====================================================================
%% The block marked "DO NOT CHANGE" must match the preamble of
%% AnonymousSubmission2027.tex in your AuthorKit exactly. If anything
%% differs (e.g. the \pdfinfo TemplateVersion), copy it from there.
%% =====================================================================

\documentclass[letterpaper]{article} % DO NOT CHANGE THIS
\usepackage[submission]{aaai2027}  % DO NOT CHANGE THIS
% times/helvet/courier are NOT loaded here: aaai2027.sty already loads
% newtxtext, helvet and courier, and forbids \usepackage{times}.
\usepackage[hyphens]{url}  % DO NOT CHANGE THIS
\usepackage{graphicx} % DO NOT CHANGE THIS
\urlstyle{rm} % DO NOT CHANGE THIS
\def\UrlFont{\rm}  % DO NOT CHANGE THIS
\usepackage{natbib}  % DO NOT CHANGE THIS AND DO NOT ADD ANY OPTIONS TO IT
\usepackage{caption} % DO NOT CHANGE THIS AND DO NOT ADD ANY OPTIONS TO IT
\frenchspacing  % DO NOT CHANGE THIS
\setlength{\pdfpagewidth}{8.5in} % DO NOT CHANGE THIS
\setlength{\pdfpageheight}{11in} % DO NOT CHANGE THIS

% Recommended by the template for listings
\usepackage{newfloat}
\usepackage{listings}
\DeclareCaptionStyle{ruled}{labelfont=normalfont,labelsep=colon,strut=off} % DO NOT CHANGE THIS
\lstset{%
	basicstyle={\footnotesize\ttfamily},
	numbers=left,numberstyle=\footnotesize,xleftmargin=2em,
	aboveskip=0pt,belowskip=0pt,
	showstringspaces=false,tabsize=2,breaklines=true}

\pdfinfo{
/TemplateVersion (2027.1)
}

% Section numbers are needed because the text uses Appendix~\ref{...}.
% The template allows changing this to 1 or 2.
\setcounter{secnumdepth}{0}

%% ---------------------------------------------------------------------
%% Your packages (only AAAI-allowed ones).
%% Removed from the NeurIPS version: neurips_2026, natbib options,
%% inputenc, fontenc (forbidden), microtype, hyperref (forbidden),
%% duplicate booktabs/array, float-fraction overrides.
%% ---------------------------------------------------------------------
\usepackage{amsmath,amssymb}
\usepackage{booktabs}
\usepackage{array}
\usepackage{tabularx}
\usepackage{xcolor}
\usepackage{tikz}
\usetikzlibrary{positioning,fit,backgrounds,arrows.meta,calc,shadows}
\usepackage{cleveref} % works without hyperref
\usepackage{xspace}

% Personal comment commands (comment out the color versions before submitting)
\newcommand{\SC}[1]{\textcolor{violet}{Simon: #1}}
\newcommand{\JB}[1]{\textcolor{teal}{Jaloliddin: #1}}
% \renewcommand{\SC}[1]{}
% \renewcommand{\JB}[1]{}

% Convenience macros
\newcommand{\rlm}{\textsc{RLM}\xspace}
\newcommand{\rlms}{\textsc{RLM}s\xspace}
\newcommand{\sh}{\textsc{Self-Harness}\xspace}
\newcommand{\sid}[1]{\textbf{\texttt{#1}}}
\newcommand{\hzero}{\ensuremath{\mathrm{H}_{0}}\xspace}
\newcommand{\hzerostar}{\ensuremath{\mathrm{H}_{0}^{*}}\xspace}

% ---- styling for the RLM edit-space figure (fig:surfaces) ----
\definecolor{editgreen}{HTML}{2E7D5B}
\definecolor{editfill}{HTML}{E7F5EE}
\definecolor{fixgrey}{HTML}{8A8F98}
\definecolor{fixfill}{HTML}{F1F2F4}
\definecolor{envtint}{HTML}{EAF1FB}
\definecolor{subtint}{HTML}{FCF1E7}
\definecolor{inkcol}{HTML}{2B2F36}
\newcommand{\cardhd}{\bfseries}
\newcommand{\cardsb}{\normalfont\scriptsize\color{black!62}}
\tikzset{
  penbadge/.pic={
    \fill[editgreen] (0,0) circle (0.095);
    \draw[white, line width=0.55pt, line cap=round] (-0.045,-0.045) -- (0.03,0.03);
    \fill[white] (-0.052,-0.052) circle (0.017);
  },
  lockbadge/.pic={
    \fill[fixgrey] (0,0) circle (0.095);
    \draw[white, line width=0.6pt] (-0.032,0.006) arc (180:0:0.032);
    \fill[white, rounded corners=0.2pt] (-0.05,-0.05) rectangle (0.05,0.014);
  },
}

% ---- code-listing style for the initial-harness figure (Figure_1.tex) ----
\definecolor{lstbg}{HTML}{FAFAFB}
\definecolor{lstframe}{HTML}{D5D8DD}
\definecolor{lstkw}{HTML}{1052B8}
\definecolor{lststr}{HTML}{B25A05}
\definecolor{lstnum}{HTML}{9AA0A8}
\definecolor{lstcom}{HTML}{6B7280}
\lstdefinestyle{harness}{
  language=Python,
  basicstyle=\ttfamily\scriptsize\linespread{0.88}\selectfont,
  keywordstyle=\color{lstkw},
  stringstyle=\color{lststr},
  commentstyle=\color{lstcom},
  deletekeywords={None,True,False},
  emph={str,list,dict,int,tuple},
  emphstyle=\color{lstkw},
  numbers=left,
  numberstyle=\ttfamily\tiny\color{lstnum},
  numbersep=9pt,
  backgroundcolor=\color{lstbg},
  frame=single,
  rulecolor=\color{lstframe},
  framesep=5pt,
  breaklines=true,
  breakindent=14pt,
  breakautoindent=false,
  columns=fullflexible,
  keepspaces=true,
  showstringspaces=false,
  xleftmargin=20pt,
  xrightmargin=1pt,
  aboveskip=2pt,
  belowskip=2pt,
}

%% ---------------------------------------------------------------------
\title{Self-Harnessing Recursive Language Models}
% Anonymous submission: keep the author block exactly as in
% AnonymousSubmission2027.tex if it differs from this.
\author{
    Anonymous Submission
}
\affiliations{
}

\begin{document}

\maketitle
% ================================================================================================
%-------------------------------------- Abstract -------------------------------------------------
% ================================================================================================

% \begin{abstract}
% Recursive language models (\rlms) use a persistent execution environment and
% recursive model calls to work over inputs larger than a single model context.
% Their performance depends on the surrounding harness. We adapt
% \sh{}~\citep{zhang2026selfharness} to this setting: a fixed model diagnoses failed
% executions, proposes bounded edits to declared harness surfaces, and compares
% one combined candidate with the incumbent on validation tasks. The current gate
% allows ties in verifier-pass count within configured resource bands. We describe
% the implemented loop and a forthcoming OOLONG-Pairs study comparing the saved
% initial and final harnesses on held-out task instances and larger record
% collections. Results and cross-environment transfer evaluation remain forthcoming.
% \end{abstract}


\begin{abstract}
\JB{Abstract from Short paper}
Recursive language models (\rlms) let a bounded-context model operate over inputs larger
than its context window by offloading the input, decomposing it, invoking recursive model
calls, and aggregating their outputs. Vanilla RLMs suffer from many failure
modes, necessitating the need to fine-tune LLMs or extend the RLM harness by hand to overcome limitations in specific tasks. We propose adapting the \sh{} framework to \rlms, so that a fixed model improves its own recursive harness from evidence produced while running: it mines recurring failure mechanisms from verifier-grounded execution traces, proposes minimal edits to declared harness surfaces, and promotes only edits that survive non-regressive
validation. The optimized harness is frozen and evaluated on instances $8$--$32\times$
longer than any seen during optimization and in new target environments. Results forthcoming.
\end{abstract}

% ================================================================================================
%------------------------------------- Introduction ----------------------------------------------
% ================================================================================================

\section{Introduction}
Frontier language models are bounded by a fixed context window, and quality degrades on
inputs longer than what the model was effectively trained and evaluated on, well before the
window limit is reached, a failure mode termed \emph{context rot}~\citep{hong2025contextrot}.
An input is called \emph{in-distribution} when its length and structure resemble what the model reliably handles, and \emph{out-of-distribution} otherwise, where behavior degrades unpredictably. Recursive language models (\rlms) scale inference-time compute instead of the
context window~\citep{snell2024scaling}: a model decomposes its own context and invokes fresh
copies of itself on the pieces~\citep{zhang2025rlm}, turning a long out-of-distribution input
into sub-calls that are individually shorter~\citep{zhang2026harness}; shortening addresses
the length dimension of the mismatch, though it does not by itself guarantee a sub-call's
content is otherwise in-distribution.
Performance therefore depends on the harness governing recursion, not just the model, and multiple possible failure modes implicate the harness:
policies collapse into a single sub-call~\citep{zhang2026harness},
deeper recursion degrades accuracy while inflating cost~\citep{wang2026overthink}, and recursion
halts by heuristic rather than evidence~\citep{guha2026termination}.

Three approaches are mainly used today.
Fine-tuning the model to work within the
harness~\citep{zhang2026harness,yang2026recursive} ties any gain to a specific checkpoint.
Hand-engineering the harness~\citep{alizadeh2026srlm,roy2026ycombinator,lumer2026rah} requires
frequent expert work that a new model family can outpace.
The third, \emph{meta-agent} approach edits the scaffold directly by inspecting an agent's own trajectories, as in \sh{}~\citep{zhang2026selfharness}, which mines execution traces for recurring failures, proposes minimal edits, and promotes only non-regressive ones. We adapt \sh{} to \rlms (SH-RLM). \rlms are a natural but untested setting for this third approach: final answers can be checked by an environment verifier, while execution traces support model-inferred diagnoses of recursive behavior.
We ask: \emph{can a language model
identify recurring failures in its own recursive behavior, convert them into targeted harness
edits, and build a harness that keeps performing when applied to inputs longer than those it was optimized on? Is a SH-RLM competitive against a hand designed harness-RLM?}

% ================================================================================================
%----------------------------------- Related Work ------------------------------------------------
% ================================================================================================

\section{Related Work}

\paragraph{Harness engineering for recursive agents.}
Recent work has explored harnesses as an important component of recursive language models (RLMs). Prime Agent extends the RLM framework with a persistent execution environment and a Continual Harness that maintains histories, memories, skills, and subagent specifications across trajectories~\citep{karten2026primeagent}. This improves the ability of a fixed RLM harness to accumulate state over time, but does not directly optimize the control logic of the harness. Our work instead studies whether targeted edits to the recursive control loop can be discovered from failure evidence and validated before deployment.

\paragraph{Self-optimizing harnesses.}
Several recent works treat the harness itself as an object of optimization. Agentic Harness Engineering uses execution traces, an explicit representation of editable components, and outcome-based validation to iteratively improve coding-agent harnesses~\citep{lin2026agenticharness}. Self-Harness similarly explores automatic harness improvement for agentic coding tasks~\citep{zhang2026selfharness}. These works demonstrate that harness changes can transfer across tasks and model families, but focus on coding agents with relatively fixed tool and action structures. Our work extends this direction to RLMs, where recursion depth and sub-call decomposition are themselves part of the editable control loop.

\paragraph{Harness optimization through model co-evolution.}
Another line of work uses a separate model or joint optimization to improve harnesses. Harness-R1 trains a dedicated harness engineer using reinforcement learning to generate executable harness patches from agent failures~\citep{shao2026harnessr1}. HarnessForge and HarnessX further investigate joint harness and policy evolution, using failure signals and structured harness components to guide optimization~\citep{chen2026harnessforge,chen2026harnessx}. In contrast, our approach uses the same frozen model for proposing and executing harness changes, with no weight updates or auxiliary trained engineer. This isolates the contribution of harness optimization itself.

\paragraph{Evaluating harness evolution.}
Recent work has also highlighted the difficulty of evaluating automatically evolved harnesses. Rethinking the Evaluation of Harness Evolution argues that improvements should be compared against task-level search under matched inference and feedback budgets, since iterative harness evolution may otherwise benefit from additional search~\citep{wang2026rethinkeval}. Harness-Bench provides a shared benchmark for evaluating harness effects across models and execution environments~\citep{yao2026harnessbench}. Our evaluation addresses these concerns by freezing the optimized harness before testing and evaluating on longer, held-out instances and a disjoint target environment.

\paragraph{Harness benefit and model capability.}
Finally, Disentangling Harness Evolution separates harness-updating from harness-benefit: models may differ in their ability to produce useful harness updates and in their ability to correctly use those updates~\citep{lin2026disentangling}. This is relevant to our setting because both harness proposal and execution use a single frozen backbone. Consequently, performance may depend not only on the quality of the discovered harness, but also on the model's ability to follow and benefit from its changes.

\section{Our Contributions}

Our work makes the following contributions:

\begin{itemize}
    \item \textbf{\sh{} for recursive language models.}
    Adapting \sh{} \citep{zhang2026selfharness} to \rlms by declaring ten editable harness surfaces keyed to phases of the \rlm{} loop, starting from a deliberately sparse initial harness and letting a fixed model mine its own traces, propose targeted edits, and validate them without weight updates or a stronger external model.

    \item \textbf{Failure-driven harness optimization.}
    We introduce a propose--validate--promote loop that uses execution failures to identify weaknesses in the harness, generate targeted modifications, and retain only validated improvements.

    \item \textbf{Length generalization.}
    We evaluate whether harnesses optimized on shorter contexts continue to improve performance on substantially longer inputs, testing whether the learned harness improvements generalize beyond the context lengths used during optimization.

    \item \textbf{Empirical evaluation.}
    We evaluate self-harnessing across long-context benchmarks and compare optimized harnesses against fixed-harness baselines to measure the effectiveness of automatically discovered harness improvements.
\end{itemize}

% \section{Our Contributions}
% \JB{Contributions from short paper. Are we not planning to do cross-environment transfer?}
% \begin{enumerate}
% \setlength{\itemsep}{1pt}
% \setlength{\topsep}{2pt}
% \setlength{\parskip}{0pt}
% \item \textbf{\sh{} for recursive language models.}
% Adapting \sh{} \citep{zhang2026selfharness} to \rlms by declaring ten editable harness surfaces keyed to phases of the \rlm{} loop, starting from a deliberately sparse initial harness and letting a fixed model mine its own traces, propose targeted edits, and validate them without weight updates or a stronger external model.

% \item \textbf{Evidence of length generalization.}
% Optimising only on short source instances, freeze the harness, and evaluate on source instances $8$--$32\times$ longer.

% \item \textbf{Evidence of cross-environment transfer.}
% Applying the same frozen harness to an unseen target environment's short and long splits.
% \end{enumerate}


% Paths are prefixed with aaai27/ because Overleaf compiles from the project
% root, where the NeurIPS files with the same names live.
% (\textwidth -> \linewidth, or figure -> figure* for full-width).
% Shared main-text method for both active manuscripts.
\section{Preliminaries and Self-Harness Optimization}
\label{sec:method}

An \rlm{} drives a persistent Python REPL to inspect an input, issue model calls,
and assemble an answer~\citep{zhang2025rlm}. The input is initially stored in
\texttt{context}; code can expose selected content through printed output. The
ordinary submission path sets \texttt{answer["content"]} and
\texttt{answer["ready"]}. Model weights, experiment-owned resource caps, and
the environment verifier remain fixed during optimization.

Following \citet{zhang2026selfharness}, we expose ten editable harness surfaces:
five instruction fields, runtime policy, execution-output formatting, REPL
helpers, answer middleware, and a skill library. Each surface corresponds to a
declared field or group of fields. Human-authored capability descriptions and
mechanism-to-surface routes constrain the search. The current study starts from
\texttt{H0*R}, which already contains hand-authored recursion and decomposition
guidance; the sparse \hzero{} is a reference baseline. Appendix~\ref{app:surfaces}
specifies the actual inputs, scope, and invocation conditions of each surface.


\subsection{Self-Harness Optimization}
\label{sec:sho}

The same fixed model executes tasks, diagnoses failures, and proposes edits.
Each round has three stages (Figure~\ref{fig:sh-loop}).

\paragraph{Weakness mining.}
The incumbent runs on mining tasks. The environment verifier evaluates final
answers, and an attributor uses selected execution-trace operations to infer
failure mechanisms with explicit uncertainty. Diagnoses use a predefined
vocabulary and are clustered by their recorded failure signature. Cited
operations support an interpretation; they do not independently establish its
causal correctness.

\paragraph{Harness proposal.}
The proposer receives compact operation-level evidence, the incumbent surfaces,
runtime capabilities, and bounded history of attempted edits and outcomes. It
selects an intervention before writing its replacement text and must explain
what the incumbent does, what observed behavior remains wrong, and what changes
at the relevant operation. At most one admitted edit may occupy a surface or
target a pattern in a round. The edit ceiling permits fewer edits or abstention.
Host materialization and preflight reject malformed or structurally invalid
changes before task evaluation.

\paragraph{Joint validation and promotion.}
All admitted edits are composed into one candidate, which is compared with a
freshly evaluated incumbent on the same validation tasks. Validation tasks are
excluded from mining but reused for adaptive selection; aggregate outcomes
inform later proposal history. In the current configuration, promotion requires
at least as many environment-verifier passes as the incumbent and compliance
with the configured resource bands. Optional verifier-defined dense metrics
inform diagnosis and history, but are not promotion conditions. The combined
candidate is the measured unit; its outcome does not identify each edit's
individual contribution. Optimization stops at the round ceiling or after the
configured number of rounds without promotion, then freezes the final incumbent
for separate evaluation. Appendices~\ref{app:mining}--\ref{app:measurement}
give the implementation contracts and measurement definitions.

% Figure 2: one Self-Harness optimization round as a loop, with pictogram panels.
% Relies on the palette (editgreen, editfill, fixgrey, fixfill, envtint, subtint,
% inkcol), card macros (\cardhd, \cardsb), and the penbadge/lockbadge pics defined
% in the preamble of proposal.tex.
\definecolor{rejred}{HTML}{B3363D}
\definecolor{minehd}{HTML}{2B5FAD}
\definecolor{valhd}{HTML}{B26A05}
%
% ---- glyph library (figure-local) -------------------------------------------
% Nested tikzpictures inherit nothing from the outer picture; glyphrow restores
% fonts, colors, and arrow tips for the in-box panels.
\tikzset{
  % sharp corners: the enclosing stage/harnessbox node's rounded corners=7pt
  % stays in force inside the node text's nested tikzpicture and would turn
  % every cornered path (checkmarks, cells) into a smooth swoosh.
  glyphrow/.style={font=\sffamily\scriptsize, line width=0.6pt, inkcol,
    sharp corners, >={Stealth[round,length=3pt]}, every node/.style={text=inkcol,
    inner sep=1pt, align=center}, baseline=(current bounding box.center)},
  % task card: a mining-data task instance (plain document)
  taskcard/.pic={
    \draw[rounded corners=1.5pt, fill=white, draw=black!40] (-0.42,-0.5) rectangle (0.42,0.5);
    \foreach \y in {0.3,0.1,-0.1,-0.3}{
      \draw[line width=0.5pt, black!40] (-0.26,\y) -- (0.28,\y);}
  },
  % Execution call tree; uniform nodes carry no semantic correctness labels.
  calltree/.pic={
    % edges first, dots after, so the dots sit on top
    \draw[black!55] (0,0.55) -- (-0.45,0.0);
    \draw[black!55] (0,0.55) -- (0,0.0);
    \draw[black!55] (0,0.55) -- (0.45,0.0);
    \draw[black!55] (0.45,0.0) -- (0.27,-0.48);
    \draw[black!55] (0.45,0.0) -- (0.63,-0.48);
    \fill[inkcol] (0,0.55) circle (0.07);
    \fill[inkcol] (-0.45,0.0) circle (0.07);
    \fill[inkcol] (0,0.0) circle (0.07);
    \fill[inkcol] (0.45,0.0) circle (0.07);
    \fill[inkcol] (0.27,-0.48) circle (0.07);
    \fill[inkcol] (0.63,-0.48) circle (0.07);
  },
  % two failure clusters grouped by signature
  clusterpair/.pic={
    \fill[rejred] (-0.38,0.1) circle (0.07);
    \fill[rejred] (-0.52,-0.12) circle (0.07);
    \fill[rejred] (-0.24,-0.14) circle (0.07);
    \draw[rejred, dashed, line width=0.5pt] (-0.38,-0.04) ellipse (0.3 and 0.3);
    \fill[valhd] (0.38,0.12) circle (0.07);
    \fill[valhd] (0.24,-0.12) circle (0.07);
    \fill[valhd] (0.52,-0.1) circle (0.07);
    \draw[valhd, dashed, line width=0.5pt] (0.38,-0.02) ellipse (0.3 and 0.3);
  },
  % the fixed model acting as proposer; the H_t subtitle (a node placed in the
  % panel, not here — nodes inside pics can drift off the pic's paths) marks
  % the current harness it edits
  proposercard/.pic={
    \draw[rounded corners=2pt, fill=white, draw=black!40] (-0.45,-0.68) rectangle (0.45,0.55);
    \fill[inkcol] (0,0.22) circle (0.11);
    \fill[inkcol] (-0.2,-0.18) arc (180:0:0.2) -- cycle;
    \path (0.45,0.55) pic[scale=1.3] {lockbadge};
  },
  % candidate-edit mark: same size as the acceptreject accept mark, which
  % renders cleanly at figure scale (the tiny penbadge strokes do not)
  checkbadge/.pic={
    \fill[editgreen] (0,0) circle (0.16);
    \draw[white, line width=1.1pt, line cap=round] (-0.08,0) -- (-0.02,-0.08) -- (0.1,0.1);
  },
  % Incumbent and combined candidate on the same held-out validation tasks.
  splitcheck/.pic={
    \draw[rounded corners=1.5pt, fill=white, draw=black!40] (-0.45,0.06) rectangle (0.45,0.62);
    \draw[line width=0.5pt, black!40] (-0.3,0.44) -- (0.18,0.44);
    \draw[line width=0.5pt, black!40] (-0.3,0.24) -- (0.28,0.24);
    \draw[editgreen, line width=0.8pt, line cap=round] (0.2,0.4) -- (0.27,0.33) -- (0.4,0.5);
    \draw[rounded corners=1.5pt, fill=white, draw=black!50] (-0.45,-0.62) rectangle (0.45,-0.06);
    \draw[line width=0.5pt, black!40] (-0.3,-0.24) -- (0.18,-0.24);
    \draw[line width=0.5pt, black!40] (-0.3,-0.44) -- (0.28,-0.44);
    \draw[editgreen, line width=0.8pt, line cap=round] (0.2,-0.28) -- (0.27,-0.35) -- (0.4,-0.18);
  },
  % Verifier-pass counts and configured resource bands.
  bandgauge/.pic={
    \draw[black!45, line width=1.1pt] (-0.45,0) arc (180:0:0.45);
    \draw[editgreen, line width=2pt] (-0.32,0.32) arc (135:45:0.45);
    \draw[inkcol, line width=0.8pt] (0,0) -- (75:0.34);
    \fill[inkcol] (0,0) circle (0.05);
    \draw[black!45, line width=0.6pt] (-0.45,0) -- (0.45,0);
  },
  % promotion decision marks
  acceptreject/.pic={
    \fill[editgreen] (0,0.24) circle (0.16);
    \draw[white, line width=1.1pt, line cap=round] (-0.08,0.24) -- (-0.02,0.16) -- (0.1,0.34);
    \fill[rejred] (0,-0.24) circle (0.16);
    \draw[white, line width=1.1pt, line cap=round] (-0.08,-0.16) -- (0.08,-0.32);
    \draw[white, line width=1.1pt, line cap=round] (-0.08,-0.32) -- (0.08,-0.16);
  },
  % Surface-grid cells (S1-S9, reading order). Each cell is ONE node whose own
  % border is the cell frame and whose text is the label — separate \draw
  % rectangles + \node labels drift apart in this document (path-vs-node offset
  % triggered by earlier nested pictures), and a single node cannot misalign
  % with itself.
  cellnode/.style={rounded corners=1pt, inner sep=0pt, minimum width=0.38cm,
    minimum height=0.38cm, font=\sffamily\tiny, align=center},
  cellplain/.style={cellnode, fill=white, draw=fixgrey, line width=0.5pt, text=black!50},
  celledited/.style={cellnode, fill=editgreen!22, draw=editgreen, line width=0.9pt, text=editgreen},
  cellnew/.style={cellnode, fill=editgreen, draw=editgreen, line width=0.9pt, text=white},
}
%
% AAAI: the three-column loop diagram spans both columns (figure*, [t] only).
\begin{figure*}[t]
\centering
\resizebox{\textwidth}{!}{%
\begin{tikzpicture}[
  font=\sffamily\small, line width=0.7pt, inkcol,
  >={Stealth[round,length=5pt]},
  every node/.style={text=inkcol},
  stage/.style={rounded corners=7pt, draw=black!30, line width=1pt, align=center,
    inner sep=8pt, text width=64mm,
    drop shadow={shadow xshift=0.4pt, shadow yshift=-0.7pt, opacity=0.22, fill=black!60}},
  harnessbox/.style={rounded corners=7pt, draw=black!30, line width=1pt, align=center,
    inner sep=8pt, text width=48mm,
    drop shadow={shadow xshift=0.4pt, shadow yshift=-0.7pt, opacity=0.22, fill=black!60}},
  flow/.style={->, black!55, line width=1.4pt},
  flowlab/.style={font=\sffamily\scriptsize\itshape, text=black!55},
]

% ================= middle column: the three stages =================
\node[stage, fill=envtint] (mine)
  {{\cardhd\color{minehd} Weakness Mining}\\[4pt]
   \begin{tikzpicture}[glyphrow]
     \pic at (0,0) {taskcard};
     \node at (0,-0.95) {run $\mathrm{H}_t$\\ mining data};
     \draw[->, black!55] (0.6,0) -- (1.0,0);
     \pic at (1.7,0) {calltree};
     \node at (1.7,-0.95) {trace-based\\ diagnosis};
     \draw[->, black!55] (2.45,0) -- (2.85,0);
     \pic at (3.65,0) {clusterpair};
     \node at (3.65,-0.95) {cluster\\ by $\phi$};
   \end{tikzpicture}};

\node[stage, fill=editfill, below=10mm of mine] (prop)
  {{\cardhd\color{editgreen} Harness Proposal}\\[4pt]
   \begin{tikzpicture}[glyphrow]
     % stacked proposer inputs, each with its own arrow into the proposer card;
     % arrows are anchored on the text nodes so they land at the right height
     % (the two-line item's arrow starts at its vertical center)
     \node[anchor=east, align=right, font=\sffamily\tiny] (inp1) at (0.35,0.52) {operation evidence};
     \node[anchor=east, align=right, font=\sffamily\tiny] (inp2) at (0.35,0.1) {surface\\ capabilities};
     \node[anchor=east, align=right, font=\sffamily\tiny] (inp3) at (0.35,-0.32) {prior edits};
     \foreach \n in {inp1,inp2,inp3}{
       \draw[->, black!55] ([xshift=2pt]\n.east) -- ([xshift=17pt]\n.east);}
     \pic at (1.55,0.1) {proposercard};
     \node[font=\tiny, text=black!55] at (1.55,-0.27) {$\mathrm{H}_t$};
     \node at (1.55,-0.8) {proposer\\(fixed model)};
     % same S1-S9 grid as the harness boxes; its middle row sits at the proposer
     % card's vertical center (0.035), and the arrow is anchored on the middle-
     % row cell so all three land on one line
     \foreach \j [count=\r from 0] in {2,1,0}{
       \foreach \i in {0,1,2}{
         \pgfmathtruncatemacro{\slab}{\r*3+\i+1}
         \ifnum\slab=4 \tikzset{c/.style=celledited}\else
         \ifnum\slab=8 \tikzset{c/.style=celledited}\else
           \tikzset{c/.style=cellplain}\fi\fi
         \node[c] (propcell\slab) at ({3.1+\i*0.46},{\j*0.46-0.15}) {S\slab};}}
     \draw[->, black!55] ([xshift=-20pt]propcell4.west) -- ([xshift=-9pt]propcell4.west);
     \path (propcell4.north east) pic {checkbadge};
     \node[cellplain] (propcell10) at ({3.1+1*0.46},{-1*0.46-0.15}) {S10};
     \node at (3.56,-1.05) {at most one edit\\ per surface};
   \end{tikzpicture}};

\node[stage, fill=subtint, below=10mm of prop] (val)
  {{\cardhd\color{valhd} Proposal Validation}\\[4pt]
   \begin{tikzpicture}[glyphrow]
     \pic at (0,0) {splitcheck};
     \node at (0,-1.05) {paired runs\\ held-out only};
     \draw[->, black!55] (0.65,0) -- (1.05,0);
     \pic at (1.8,-0.15) {bandgauge};
     \node at (1.8,-1.05) {pass counts +\\ resource bands};
     \draw[->, black!55] (2.55,0) -- (2.95,0);
     \pic at (3.5,0) {acceptreject};
     \node at (3.5,-1.05) {promote /\\ reject batch};
   \end{tikzpicture}};

% ================= side columns: current and updated harness =================
\node[harnessbox, fill=fixfill, left=8mm of prop] (cur)
  {{\cardhd Current RLM Harness $\mathrm{H}_t$~\tikz[baseline=-2pt, scale=0.8,
     sharp corners]{% sharp corners: the box node's rounded corners=7pt leaks
     % into this inline picture and would distort the small glyph paths
     \fill[inkcol] (0,0.22) circle (0.11);
     \fill[inkcol] (-0.2,-0.18) arc (180:0:0.2) -- cycle;}}\\[4pt]
   \begin{tikzpicture}[glyphrow]
     % all placement via node anchors: raw-coordinate paths drift relative to
     % nodes in this document, but paths anchored on node coordinates land true
     \node[rounded corners=2pt, draw=fixgrey, dashed, line width=0.7pt, fill=white,
       minimum width=1.06cm, minimum height=1.0cm, align=center] (mc) at (0,0) {fixed\\ model};
     \path (mc.north east) pic[scale=1.3] {lockbadge};
     \foreach \j [count=\r from 0] in {2,1,0}{
       \foreach \i in {0,1,2}{
         \pgfmathtruncatemacro{\slab}{\r*3+\i+1}
         \ifnum\slab=2 \tikzset{c/.style=celledited}\else
         \ifnum\slab=6 \tikzset{c/.style=celledited}\else
           \tikzset{c/.style=cellplain}\fi\fi
         \node[c] (curcell\slab) at ({1.62+\i*0.46},{\j*0.46-0.46}) {S\slab};}}
     % + drawn as path strokes anchored on node coordinates (a text node here
     % drifts off the true midpoint)
     \coordinate (plusmid) at ($(mc.east)!0.5!(curcell4.west)$);
     \draw[black!55, line width=1.1pt, line cap=round]
       ($(plusmid)+(-0.12,0)$) -- ($(plusmid)+(0.12,0)$);
     \draw[black!55, line width=1.1pt, line cap=round]
       ($(plusmid)+(0,-0.12)$) -- ($(plusmid)+(0,0.12)$);
     \node[cellplain] (curcell10) at ({1.62+1*0.46},{-1*0.46-0.46}) {S10};
     \node[anchor=north] (eslab) at ([yshift=-4.5pt]curcell10.south) {Editable Surfaces};
     \node[anchor=north, font=\sffamily\tiny, text=black!55] at ([yshift=-1pt]eslab.south) {green: prev. edited};
   \end{tikzpicture}};

\node[harnessbox, fill=editfill, right=8mm of prop] (new)
  {{\cardhd Updated RLM Harness $\mathrm{H}_{t+1}$~\tikz[baseline=-2pt, scale=0.8,
     sharp corners]{% sharp corners: see the current-harness title note
     \fill[inkcol] (0,0.22) circle (0.11);
     \fill[inkcol] (-0.2,-0.18) arc (180:0:0.2) -- cycle;}}\\[4pt]
   \begin{tikzpicture}[glyphrow]
     \node[rounded corners=2pt, draw=fixgrey, dashed, line width=0.7pt, fill=white,
       minimum width=1.06cm, minimum height=1.0cm, align=center] (mcn) at (0,0) {fixed\\ model};
     \path (mcn.north east) pic[scale=1.3] {lockbadge};
     \foreach \j [count=\r from 0] in {2,1,0}{
       \foreach \i in {0,1,2}{
         \pgfmathtruncatemacro{\slab}{\r*3+\i+1}
         \ifnum\slab=2 \tikzset{c/.style=celledited}\else
         \ifnum\slab=4 \tikzset{c/.style=cellnew}\else
         \ifnum\slab=8 \tikzset{c/.style=cellnew}\else
         \ifnum\slab=6 \tikzset{c/.style=celledited}\else
           \tikzset{c/.style=cellplain}\fi\fi\fi\fi
         \node[c] (newcell\slab) at ({1.62+\i*0.46},{\j*0.46-0.46}) {S\slab};}}
     \coordinate (plusmidn) at ($(mcn.east)!0.5!(newcell4.west)$);
     \draw[black!55, line width=1.1pt, line cap=round]
       ($(plusmidn)+(-0.12,0)$) -- ($(plusmidn)+(0.12,0)$);
     \draw[black!55, line width=1.1pt, line cap=round]
       ($(plusmidn)+(0,-0.12)$) -- ($(plusmidn)+(0,0.12)$);
     \node[cellplain] (newcell10) at ({1.62+1*0.46},{-1*0.46-0.46}) {S10};
   \end{tikzpicture}\\[4pt]
   {\cardsb joint update at \sid{S4} and \sid{S8}}};

% ================= loop arrows =================
% current harness up into weakness mining
\draw[flow] (cur.north) |- node[flowlab, pos=0.72, above]{run on tasks} (mine.west);
% mining -> proposal -> validation
\draw[flow] (mine.south) -- node[flowlab, right=1mm]{failure patterns} (prop.north);
\draw[flow] (prop.south) -- node[flowlab, right=1mm, align=left]{admit + compose\\ one candidate} (val.north);
% all rejected: back to the same harness
\draw[flow, rejred] (val.west) -| node[flowlab, text=rejred, pos=0.28, below]{no promotion --- keep $\mathrm{H}_t$} (cur.south);
% Promotion adopts the already composed and measured candidate.
\draw[flow, editgreen] (val.east) -| node[flowlab, text=editgreen, pos=0.28, below]{promote measured batch} (new.south);
% updated harness starts the next round
\draw[flow, dashed] (new.north) -- ++(0,12.5mm) node[flowlab, above=0.5mm, align=center]{next round\\ ($\mathrm{H}_{t+1}$ becomes current)};

\end{tikzpicture}}
\caption{One \sh{} round. Final-answer outcomes and selected execution operations
support model-inferred diagnoses. The proposer selects interventions and supplies
edits on distinct surfaces. Admitted edits are composed \emph{before} comparing
one candidate with a fresh incumbent on the same held-out validation tasks.
The current gate permits ties in verifier-pass count within resource bands;
the whole batch is promoted or rejected. Validation summaries inform later
history, but validation traces are not mined. An empty or skipped batch schedules
no validation. Side grids illustrate prior edits and two jointly promoted edits;
they are not individual efficacy measurements. Weights and evaluator stay fixed.}
\label{fig:sh-loop}
\end{figure*}
% Shared main-text experimental protocol; results are prospective.
\section{Experiments}
\label{sec:setup}

The forthcoming study uses the current OOLONG-Pairs profile. Runner, attributor,
and proposer use the same frozen \textbf{DeepSeek-V4-Flash} deployment and
decoding settings. Mining, proposal, validation, checkpointing, and final
evaluation are implemented; results for this study remain forthcoming.

\paragraph{Data and splits.}
\label{sec:envs-and-splits}
OOLONG-Pairs is reconstructed from the \texttt{trec\_coarse} portion of
OOLONG-synth~\citep{bertsch2025oolong}, using twenty predicates over each record
collection and a deterministic scorer for sets of \emph{user-ID pairs}.
Table~\ref{tab:study-settings} gives the configured allocation. Mining,
validation, and short test contain disjoint task instances on shared
short-context records. Short-test results therefore measure performance on
held-out task instances. The separate long evaluation assesses performance on
larger record collections, without claiming entirely unseen underlying question
texts. The current split is retained. Appendix~\ref{app:protocol} specifies the
dataset pin, reconstruction, and overlap.

\begin{table}[t]
\centering\small
\begin{tabularx}{\linewidth}{@{}lrl>{\raggedright\arraybackslash}X@{}}
\toprule
Role & Tasks & Attempts/task & Use \\
\midrule
Short mining & 20 & 2 & Failure evidence \\
Short validation & 10 & 1 per harness & Joint promotion decision \\
Short test & 10 & 3 per condition & Held-out task instances \\
Long test & 40 & 3 per condition & Larger record collections \\
\bottomrule
\end{tabularx}
\caption{Configured study. At most four admitted edits per round, three
optimization rounds, and patience three. Upstream short/long window labels are
8,192/262,144 tokens; these are not
measured lengths of the reconstructed prompts.}
\label{tab:study-settings}
\end{table}

\paragraph{Comparison and metrics.}
\label{sec:baselines}
The primary comparison is the exact saved starting harness \texttt{H0*R}
(evaluation condition \texttt{initial}) against the frozen edited harness
(\texttt{sh\_rlm}), on matching short and long task sets. Reference conditions
are the sparse \hzero{} (\texttt{b1}), the upstream-style \hzerostar{}
(\texttt{h0\_star}), and a pinned \emph{paper reconstruction} of
$\lambda$-RLM~\citep{roy2026ycombinator}. The evaluation condition list must be
selected explicitly: the runner's default reference grid omits \texttt{initial}.
Exact pair-set equality defines a pass. We additionally report recorded
pair-level precision, recall, F1, and available missing/extra-pair counts,
together with usage, cost, runtime, and completion coverage. These supplementary
metrics do not alter the promotion gate. Final evaluation uses three attempts
per task; individual model attempts have no controlled random seed.

\paragraph{Interpretation and planned extensions.}
Results will compare matching tasks with completed, failed, and missing attempt
counts reported separately. Repeated attempts and predicates sharing a context
are not independent corpus samples. A promotion under one-attempt validation
is a selection decision, not a significance test. Cross-environment evaluation
and weight-training comparisons remain future work: this profile materializes
only OOLONG-Pairs short and long splits. No completed transfer result or
long-context performance gain is claimed here.

\section{Limitations and Future Work}
\label{sec:limitations}

Diagnoses and behavioral predictions are model judgments constrained by observed
operations, not causal proofs. Joint promotion measures a batch, and an
exact-pass tie can promote despite a decrease in an optional dense metric.
Validation is reused adaptively, and the short roles share record collections.
The study starts from hand-authored H0*R guidance and uses only two contexts per
length. These choices define the scope of the evaluation; stronger claims about
individual edits or unseen source data would need additional evidence.

All optimization stages and evaluation infrastructure are implemented. The next
step is to freeze the study configuration, run optimization, and evaluate the
saved initial and final harnesses under the declared protocol. Completed results,
resource totals, and any uncertainty analysis will be reported with their actual
coverage. Cross-environment transfer, weight training, and optional ablations
remain separate future studies (Appendix~\ref{app:extensions}).

\paragraph{Broader impacts.}
Harness optimization could reduce manual engineering and improve the use of
fixed models, but a benchmark pass does not measure safe deployment or resistance
to misuse. Automatically edited instructions may change behavior outside the
benchmark; the optional safety analysis in Appendix~\ref{app:alignment} is not
part of the current promotion gate. This study uses public benchmark data and
no human-subject experiment or deployed decision system.

% AAAI sets the bibliography style itself (aaai2027.bst) — no \bibliographystyle.
\bibliography{aaai27/references}

% Check the call for papers: AAAI may require the appendix as a separate
% supplementary file instead of after the references.
\clearpage
\appendix
% Shared implementation appendix. Both active papers include this file in full.
\section{Runtime and Initial Harness}
\label{app:surfaces}

The implementation described here is the repository at revision
\texttt{070d1a8493d7}, with the OOLONG-Pairs configuration inspected on
September 28, 2026. These are implementation specifications for the forthcoming
study, not experimental results. The study's eventual launch revision, frozen
configuration, split manifest, and initial/final harness hashes must accompany
its results.

\paragraph{Initialization.}
\hzero{} supplies a minimal REPL contract and generic instruction strings, a
disabled runtime policy, the default metadata formatter, empty helper and skill
libraries, and accepting answer middleware. \hzerostar{} replaces S1 with the
upstream-style system prompt, leaves S2--S5 empty, and enables the orchestrator
addendum. \texttt{H0*R}, the configured initial harness, additionally replaces
S1's recursive-call description and supplies a hand-authored S2 example. That
example distinguishes a bare completion from a child with its own REPL, splits
line-oriented input into batches of 300, requests label counts, merges counts in
Python, and retries a chunk after a parse failure. It was informed by development
observations; gains over \hzero{} alone would conflate this initialization with
automated editing. The \texttt{initial} evaluation condition loads the exact
saved round-one harness rather than substituting a reference from today's
registry. Figure~\ref{fig:edit-surfaces} and Figure~\ref{fig:full-repl-contract}
show the sparse reference, not the complete \texttt{H0*R} prompt.

\paragraph{Execution contract.}
S1--S5 are concatenated into one system prompt; a clause referring to early turns
is guidance, not a separately timed host hook. Recursive \texttt{rlm\_query}
children below the depth limit receive their own REPL and the shared prompt.
Bare \texttt{llm\_query} calls and leaves at the recursion limit are completions,
without that REPL. Printed code output joins accumulated message history;
input-as-variable does not prevent code from printing input records. The normal
answer path uses the REPL answer dictionary, while iteration exhaustion can
invoke the runtime's final-answer fallback. Admission checks preserve specified
interfaces and caps, not a formal proof of arbitrary generated code's behavior.

% AAAI: full-width table* for the long per-surface descriptions.
\begin{table*}[t]
\centering\small
\renewcommand{\arraystretch}{1.12}
\begin{tabularx}{\textwidth}{@{}p{0.17\textwidth}X@{}}
\toprule
Surface & Editable payload, visible information, and execution scope \\
\midrule
S1: REPL contract & Full text replacement describing the environment/API; shared root/recursive-child prompt. \\
S2: Decomposition & Full text replacement guiding input inspection and decomposition; shared prompt, scoped to the current subtask. \\
S3: Execution & Full text replacement guiding parsing, calls, aggregation, and task predicates; shared prompt. \\
S4: Verification & Full text replacement guiding checks before submission; shared prompt, with no independent semantic oracle. \\
S5: Recovery & Full text replacement guiding actions after visible errors; cannot extend a terminated run or add host retries. \\
S6: Runtime policy & Dictionary with enabled flag, prompt/batch caps, depth, syntax-retry/count, and output-validation settings. Enforcement is root local except inherited depth. Width/length caps refuse work; syntax retry repeats the same prompt. \\
S7: Metadata & Root formatter function receiving bounded stdout and redacted variable types/lengths; preserves its declared output bound. It cannot inspect hidden variable values or the full call history. \\
S8: REPL helpers & One named function added/replaced in one root or child helper dictionary. Helpers execute only when called; the proposer must identify a caller and input/output contract. \\
S9: Answer middleware & Root function receiving answer text and redacted inventory; may accept transformed text or redirect with a nudge. It cannot inspect hidden records or classifications. \\
S10: Skills & One named text procedure added, replaced, or removed. A shared name/description index advertises it; its body matters only when loaded or forwarded. \\
\bottomrule
\end{tabularx}
\caption{Ten surfaces and their actual capabilities. Provider/model selection,
decoding/reasoning settings, verifiers, and experiment-owned caps are outside the
edit space. S8 spans two fields; a surface is not necessarily one builder.}
\label{tab:surfaces}
\end{table*}

\paragraph{Skills and token accounting.}
The S10 index enters the system prompt after S5. For a nonempty library the
runner installs \texttt{load\_skill(name)} in root and recursive-child REPLs.
The host merges a single-entry edit with the incumbent library, including
explicit removal. Limits are eight entries, 40 characters per name, 200 per
description, 4,000 per body, and 16,000 total. A loaded body stored only in a
variable need not be model visible. If printed or forwarded it incurs input
tokens, potentially on later turns while it remains in history. Neither loading
a body nor advertising it proves that its procedure was followed.

% Figure 1: the sparse reference harness H0 and its ten editable surfaces.
% AAAI: full-width figure* (code lines are too wide for one column); figure*
% only supports [t] placement in AAAI (stfloats is forbidden).
% Body only; the `harness' listing style, its colours, and \hzero live in the
% preamble of proposal.tex, matching how Figure_2.tex relies on the tikz styling
% defined there.
\begin{figure*}[t]
    \centering
    \begin{lstlisting}[style=harness, gobble=4]
    # shrlm/rlm_harness.py: the ten surfaces Self-Harness may edit (S1 to S10).
    
    # Condensed for visibility; the full literal appears in the appendix.
    MINIMAL_REPL_CONTRACT = """
    You are an RLM operating through a persistent Python REPL.
    Use `context`, flat or recursive LM calls, `SHOW_VARS()`, and `answer`.
    Execute one ```repl``` block per turn; print inspections and mark the
    answer ready only when complete. [Display-only condensation.]
    """
    
    def build_repl_contract() -> str:                                    # S1
        return MINIMAL_REPL_CONTRACT
    
    def build_decomposition_instruction() -> str:                        # S2
        return ("Start by probing `context` to understand what you have, "
                "then decide how the task breaks into steps.")
    
    def build_execution_instruction() -> str:                            # S3
        return ("On each turn, run one block of code, print what you need "
                "to see, and build the result up in REPL variables.")
    
    def build_verification_instruction() -> str:                         # S4
        return 'Check your candidate answer before setting `answer["ready"] = True`.'
    
    def build_recovery_instruction() -> str:                             # S5
        return ("If a sub-call errors or returns something you cannot use, "
                "adjust the approach and try again.")
    
    def build_runtime_policy() -> dict[str, Any]:                        # S6
        return {
            "enabled": False,
            "max_prompt_chars": None,
            "max_batch_width": None,
            "max_depth": None,
            "retry_on_syntax_error": None,
            "max_retries": None,
            "validate_sub_output": None,
        }
    
    def build_metadata(                                                  # S7
        stdout: str,
        repl_inventory: dict[str, tuple[str, int]],
        max_character_length: int = DEFAULT_MAX_CHARACTER_LENGTH,
    ) -> str:
        return default_metadata_builder(stdout, repl_inventory,
                                        max_character_length)
    build_metadata.declared_bound = DEFAULT_MAX_CHARACTER_LENGTH        # S7
    
    def build_repl_helpers() -> dict[str, Any]:                          # S8
        return {}
    
    def build_sub_repl_helpers() -> dict[str, Any]:                      # S8
        return {}
    
    def build_answer_middleware() -> AnswerMiddlewareFn:                 # S9
        return accept_answer
    
    @dataclass(frozen=True)
    class SkillEntry:                                                    # S10
        name: str          # REPL-safe identifier; index label
        description: str   # one line: when to consult it; enters the prompt
        body: str          # ordered steps; served verbatim by `load_skill(name)`
    
    def build_skills() -> list[SkillEntry]:                              # S10
        return []
    \end{lstlisting}
    \caption{Sparse reference harness \hzero{} and its ten editable surfaces,
    each a declared field or group of fields (S8 has two builders). The configured
    study starts from \texttt{H0*R}, described in Appendix~\ref{app:surfaces}.
    S1 is condensed for display; Figure~\ref{fig:full-repl-contract} gives the
    full \hzero{} S1 template. S10 advertises a name/description index and loads
    bodies on demand; token cost depends on exposure in conversation history or
    forwarding to children. Each admitted edit changes one surface, and a round
    may compose several edits on distinct surfaces into one measured candidate.}
    \label{fig:edit-surfaces}
\end{figure*}
% Appendix figure: exact, unabridged sparse H0 S1 template (not H0*R).
% Figure 1 intentionally contains a display-only condensation of this prompt.
% AAAI: full-width figure* (the prompt lines are too long for one column);
% figure* only supports [t] placement in AAAI (stfloats is forbidden).
\begin{figure*}[t]
    \centering
    \begin{lstlisting}[style=harness, gobble=4]
    MINIMAL_REPL_CONTRACT = textwrap.dedent(
        """\
        You are a Recursive Language Model (RLM): a language model whose prompt-related information lives in a Python REPL that you drive turn by turn. You will be queried turn-by-turn until you have an answer to the query.

        Available in the REPL:
        - `context`: the information related to the prompt (typically `str` or `list[str]`).
        - `llm_query(prompt: str, model: str | None = None) -> str`: a single sub-LLM completion.
        - `llm_query_batched(prompts: list[str], model=None) -> list[str]`: several sub-LLM completions in parallel; same order out as in.
        - `rlm_query(prompt, model=None)` / `rlm_query_batched(prompts, model=None)`: recursive RLM sub-calls. Fall back to `llm_query` / `llm_query_batched` when recursion is disabled.
        - `SHOW_VARS() -> str`: list every variable currently in the REPL.
        - `answer`: dict initialized to `{{"content": "", "ready": False}}`. To submit, set `answer["content"]` to the final answer and `answer["ready"] = True` inside a ```repl``` block.
        {custom_tools_section}

        Write code in ```repl``` blocks, one block per turn; the REPL persists across turns. The REPL is NOT a Jupyter cell - only `print(...)` output (stdout) is shown back to you between turns; a bare expression on the last line is silently discarded. REPL outputs over ~20K characters are truncated."""
    )
    \end{lstlisting}
    \caption{The complete sparse \hzero{} S1 source template, before placeholder
    formatting. Figure~\ref{fig:edit-surfaces} uses a display-only condensation.
    This is a reference contract, not the full initial prompt of the configured
    \texttt{H0*R} study; its upstream-style prompt, recursive-call change, S2
    example, and orchestrator addendum are specified in
    Appendix~\ref{app:surfaces}.}
    \label{fig:full-repl-contract}
\end{figure*}

\section{Mining, Diagnosis, and Evidence Selection}
\label{app:mining}

Final-answer verification supplies pass/fail and an outcome category. Eligible
failed runs are sent to a fixed-model attributor; passing runs and configured
environment-owned exclusions are not sent through failure attribution. Excluded,
unattributed, and rejected records remain visible in accounting. Each diagnosis
records an inferred mechanism, location and causal status, supporting operation
references, and uncertainty. Structural absence of recursion can be observed
directly; semantic correctness and root-versus-child causality are not implied
by final-answer scoring.

The attributor must distinguish skipped records, wrong or uncertain labels, and
incorrect aggregation or predicates. Missing answer elements alone do not prove
input loss. It must account for observed checks and repairs, and distinguish
missing evidence from evidence that an action was omitted. For
\texttt{incomplete\_coverage}, a structured basis records status, input scope,
loss observation, and counterevidence. Statuses \texttt{not\_established} and
\texttt{contradicted} normalize to \texttt{OTHER}/\texttt{UNATTRIBUTED}, rather
than authorizing a coverage edit. Host acceptance verifies schema, vocabulary,
field bounds, and resolvable references, not the truth of a causal explanation.
Attribution response attempts default to three; transport retries are separate.

\paragraph{Clustering and routing.}
Records cluster by exact agreement on
% AAAI: split over two lines so it fits the single column.
\[
\begin{aligned}
\phi(r)=(&\text{outcome cause},\ \text{recorded failing level},\\
          &\text{causal status},\ \text{mechanism}).
\end{aligned}
\]
The signature mixes observed outcomes with model-inferred fields; it is not a
verified causal label. Patterns rank first by distinct-instance support, then
actionability, then a deterministic signature tie-break. Attempt support is
recorded separately. Minimum support two is a flag, not automatic deletion.
Three representatives and a fixed actionability heuristic guide reading order:
weights are 0.4 causal status, 0.3 grounding availability, 0.2 known mechanism,
and 0.1 normalized-description homogeneity. The last term uses distinct text
strings, not semantic embeddings; the score is not a calibrated success
probability.

\begin{table}[t]
\centering\small
\begin{tabularx}{\linewidth}{@{}Xl@{}}
\toprule
Mechanism & Eligible surfaces (primary first) \\
\midrule
REPL contract misuse & S1, S9, S10, S4 \\
Whole-input subcall collapse & S2, S3, S10, S6 \\
Incomplete coverage & S2, S4, S3, S10 \\
Redundant decomposition & S2, S3, S10 \\
Misaligned unit & S2, S10 \\
Depth degradation & S3, S6, S2 \\
Insufficient recursion & S3, S2, S6 \\
Bare completion substituted for recursion & S3, S1, S10 \\
Premature termination & S4, S9, S3 \\
Skipped verification & S4, S9, S10 \\
Swallowed subcall error & S5, S3, S6 \\
Iteration budget exhaustion & S6, S3, S1, S10, S5$^{\dagger}$ \\
Unparsed child output & S7, S3, S5, S8$^{\dagger}$ \\
REPL execution fault & S8, S10, S3, S1, S5$^{\dagger}$ \\
Lossy aggregation & S3, S4, S8$^{\dagger}$ \\
Unconsulted procedure & S10, S3, S2 \\
Other & S1--S10 \\
\bottomrule
\end{tabularx}
\caption{Human-defined, many-to-many routing (taxonomy 3.3.0).
An \texttt{UNATTRIBUTED} diagnosis has no eligible route, including when its
mechanism is Other. $\dagger$ denotes additional operation-evidence requirements.}
\label{tab:routes}
\end{table}

Conditional S8 eligibility requires a complete cited operation; unparsed child
output additionally requires child-response evidence. Conditional S5 eligibility
requires an observed error and a subsequent complete operation. These are
structural opportunities to intervene, not proof of successful recovery or an
effective proposed fix. Aggregation errors route to S3/S4 or supported S8;
S9 remains limited to answer-visible properties.

\paragraph{Bounded evidence.}
The attribution digest defaults to 12,000 characters. Proposal evidence defaults
to 32,000 characters, separately from the 12,000-character history budget.
These are component budgets, not a total prompt cap. A compact pattern inventory
is supplemented by up to $\min(k,4)$ expanded patterns, preferring distinct
actionable mechanisms before additional signatures/context. Inventory-only
patterns do not authorize proposals. Evidence follows cited child calls and
relevant parsing, merging, or filtering operations, with nearby input definitions,
consumers, and corrections. Whole code operations are retained where they fit;
required operations that cannot fit are omitted with scope markers rather than
presented as complete fragments. Outputs and prose may be shortened. A passing
contrast sharing operation names can be included when available. This improves
context for comparison but guarantees neither full causal context nor balanced
surface coverage.

\section{Proposal, Materialization, and History}
\label{app:proposal}

One response under \texttt{proposal-selection/v2} selects evidence-supported
pattern/surface interventions before supplying replacement payloads; these are
not separate selection and generation calls. The first admitted candidate owns
its surface and pattern. Rejected entries do not reserve a surface. At most
$k$ edits are admitted, with fewer or zero allowed. Three short explanation
fields identify incumbent behavior, unresolved observed failure, and the change
at a specific operation, input, or trigger. Guidance asks whether the same
demonstrated failure could persist if the edit were followed perfectly. Repeating
an existing instruction more emphatically, moving it to another surface, or
repairing an already-recovered fault is insufficient without a substantive
behavioral difference. Replacing a misleading example can be the smallest
effective change.

The general prompt asks which information must survive each operation and which
task conditions must the final computation enforce. It does not prescribe an
OOLONG-specific record-ID algorithm. Environment-provided API and answer
contracts remain authoritative. These instructions guide model judgment; they
are not deterministic tests of semantic novelty or correctness.

\paragraph{Repair and preflight.}
Response/schema attempts are bounded (default maximum eight), distinct from
task validations. When failed-member repair is invoked there is one repair
response, restricted to failed original patterns. It may select another eligible
unoccupied surface or withdraw an edit while retaining valid siblings unchanged;
it may not invent new patterns. Collision repair must choose a contender or
withdraw it, still enforcing one admitted edit per surface.

For text under \texttt{literal-text/v1}, the model supplies literal instructions;
the host encodes formatting braces exactly once and restores the designated
custom-tools placeholder. S1--S5 replace full text, S6 replaces a policy
dictionary, S7/S9 supply functions with required interfaces, S8 edits one named
helper in one dictionary, and S10 edits one named library entry. Preflight checks
serialization and base hashes, an actual single-surface change, reserved
interfaces, caps, formatting, and selected smoke/answer fixtures. Environment
answer fixtures are synthetic contract examples, not validation tasks. Literal
no-op rejection does not detect every behaviorally redundant instruction.
Saved artifacts distinguish literal proposal text, encoded source/template, and
runtime rendering.

\paragraph{History and qualified progress.}
The archive retains attempts, admission/rejection reasons, batch membership,
measured outcomes, and revision links. The bounded prompt summary prioritizes
recent rounds, the incumbent's promotion, relevant measured predecessors, and an
older promising direction when it fits. Revisiting an intervention requires a
resolvable predecessor reference and a stated substantive change to the operation
or joint hypothesis. Validation contributes aggregate scores, costs, sanitized
failure summaries, and observed activation; its raw traces and gold/answer dumps
are not forwarded to the proposer.

A rejected measured candidate can be marked potentially promising when a
comparable verifier-defined dense metric improves in its declared direction.
Definitions specify identity/version, direction, aggregation, precision, and
terminal-outcome values. Comparison also requires compatible verifier/protocol/
repetition settings and matching instance-attempt sets. Missing, unsupported, or
incomparable metrics remain \texttt{not\_assessed}. The label is qualified
evidence, not promotion or statistical significance; a simultaneous exact-pass
regression remains recorded. An identical rejected batch already measured under
the same incumbent can be skipped before new validation.

\paragraph{Observed behavior.}
Instrumentation reports root S6 syntax retries and S9 redirects, and S10 loads
across the traced call tree, with observation coverage. Missing coverage is
unassessed. The intended-behavior field remains \texttt{not\_assessed}: activation
does not establish compliance or task correctness. There is no general semantic
compliance detector for S1--S5, S7, or S8.

\section{Joint Validation and Stopping}
\label{app:validation}
\label{app:algorithm}

Let $H$ be the incumbent, $C$ the composition of all admitted edits, and
$P(X)$ the number of final-answer verifier passes from harness $X$ on the same
$n_{\mathrm{ho}}v$ validation attempts. A singleton is the combined candidate;
an empty or skipped batch schedules no validation. Only $H$ and $C$ are scored;
constituents of a multi-edit batch have membership records without individual
scores. The implemented pass-count conditions are
\begin{equation}
\Delta=P(C)-P(H),\qquad
\Delta\geq-\tau_{\mathrm{reg}},\qquad
\Delta\geq\tau_{\mathrm{imp}}.
\label{eq:promotion}
\end{equation}
For each configured resource metric $q$, its mean must lie within the inclusive
baseline-relative band, $a_q\bar q(H)\leq\bar q(C)\leq b_q\bar q(H)$, with
unbounded endpoints permitted. Current thresholds are both zero; mean cost must
be between zero and three times incumbent cost, and subcall count is unconstrained.
Thus an exact-pass tie, including zero versus zero, can promote even if F1 falls.
This is the intended aggregate verifier-pass/resource rule, not a guarantee of
per-task improvement or secondary-metric non-regression.

\begin{quote}\small
\textbf{Inputs:} configured initial harness $H$, mining tasks $M$, validation
tasks $V$, repetitions $m,v$, edit ceiling $k$, round ceiling $T$, patience $p$,
and promotion settings. Initialize history and no-promotion counter $z=0$.
\begin{enumerate}
\item For each round up to $T$, execute or resume $H$ on $M$ with $m$ attempts;
record final outcomes, traces, and available observations.
\item Diagnose eligible failures, cluster records, and select bounded evidence
and history. Propose, admit, and if needed repair up to $k$ edits on distinct
patterns and surfaces.
\item Set promoted to false. If an evaluable, nonduplicate batch exists, compose
$C$ and freshly evaluate or resume $H,C$ on $V$ with $v$ attempts. An unscorable
baseline follows the experiment error path. For a complete comparable candidate,
apply Equation~\ref{eq:promotion} and resource bands; if accepted set $H=C$ and
promoted to true.
\item Persist the batch outcome and constituent membership, or the reason no
comparison occurred; update history. Set $z=0$ on promotion, otherwise $z=z+1$.
Stop when $z\geq p$.
\item Freeze the last incumbent; run separately specified final-evaluation
conditions. Do not retrospectively select a harness using test scores or F1.
\end{enumerate}
\end{quote}

\paragraph{Failures, denominators, and resume.}
Candidate-owned execution errors can become failed attempts without terminating
the whole experiment. Configuration/operational/provider errors, cancellation,
memory or I/O failures, and observation-persistence failures may still propagate.
Complete comparisons require matching planned attempts; a partial baseline cannot
anchor promotion. The final-evaluation driver can retain completed sibling
results while reporting an incomplete or over-budget condition.

\begin{table}[t]
\centering\small
\begin{tabularx}{\linewidth}{@{}p{0.27\linewidth}X@{}}
\toprule
Recorded status & Treatment \\
\midrule
Completed correct/incorrect & Enters attempt denominator; only a verifier pass counts as success. \\
Completed execution/resource/format/filter failure & Enters denominator as a failure; dense terminal values follow the verifier definition. Partial usage/trace coverage remains labeled. \\
Missing/canceled attempt & Report separately, not as a fabricated measured zero. Incomplete validation cannot silently become a complete comparison. \\
Unscorable baseline & Persist evidence and stop/raise through the experiment error path; this is not a candidate rejection. \\
Resumed/cached completion & Reuse only with matching identity/contracts; not a new independent attempt or new spend. \\
\bottomrule
\end{tabularx}
\caption{Outcome accounting for optimization and final evaluation.}
\end{table}

\section{Experimental Protocol and Provenance}
\label{app:protocol}

\paragraph{Pinned data and scoring.}
The profile is \path{configs/experiment_oolong_pairs_DeepSeekV4Flash.toml}.
The dataset is \texttt{oolongbench/oolong-synth}, upstream validation partition,
\texttt{trec\_coarse}, revision
\path{f0d59eaf0febf130664cfceb710436c8e3216b2b}. The loader selects the first
two distinct windows per configured length, with a 50,000-row scan ceiling,
then pairs each with task IDs 1--20 and shuffles the instances with seed zero.
The seed controls task partitioning, not per-attempt model randomness. Labels
are removed from model-visible records; gold user-pair sets are computed from
the labeled records and task predicates.

The scorer extracts integer pairs and compares sets. Global pair order and
duplicate pairs do not affect set equality, but order within a pair is preserved:
the lower user ID must come first. A bracketed list or newline-separated pairs
are accepted. The explicit empty marker is ``No valid pairs found.''; literal
\texttt{[]} is not an empty answer. No parseable pair or marker is a format
failure. A correctly empty prediction against an empty gold set passes and has
precision/recall/F1 equal to one. These extraction and empty-set distinctions
are adaptations of the example scorer, not an unmodified upstream verifier.

\paragraph{Task separation and shared records.}
The two short windows contain 188 records each; the two long windows contain
6,374 each. Short instances are partitioned 20/10/10, while all 40 long
instances are reserved for separate evaluation. The current split remains
unchanged. Table~\ref{tab:context-overlap} makes its unit explicit.

\begin{table}[t]
\centering\small
\setlength{\tabcolsep}{3.5pt} % AAAI: fit the single column
\begin{tabular}{@{}lrrr@{}}
\toprule
Short context (input-prefix hash) & Mining & Validation & Short test \\
\midrule
Window 9 (\texttt{5f067945b676}) & 10 & 6 & 4 \\
Window 10 (\texttt{e939371361d4}) & 10 & 4 & 6 \\
\bottomrule
\end{tabular}
\caption{Task counts per shared short record collection at partition seed zero.
Hashes exclude the task question. Distinct task instances do not imply distinct
underlying record collections.}
\label{tab:context-overlap}
\end{table}

Different predicates over the same window ask different aggregation questions
over identical records. Between the two short windows, no complete
\texttt{(user\_id, date, question\_text)} record is identical, but 43 question
texts recur with different users/dates. The long windows also have no identical
complete-record overlap with the short windows, while sharing question texts.
Short-test results support the held-out-task-instance interpretation; long
evaluation assesses larger record collections. Neither establishes wholly
unseen-source-text generalization. A stronger source-separated design is needed
only if that becomes a central claim. This overlap is not evidence of answer
memorization.

The nominal upstream lengths differ by $32\times$, but stripping labels and
adding task text changes actual prompt lengths. Evaluation will report
reconstructed prompt character counts and, when measured, token counts with
the tokenizer identified, separately from nominal window labels. Multiple
predicates per window and three requests per task do not increase the number of
independent contexts. Any future interval or paired test must specify its unit,
attempt aggregation, missingness, and the two-context limitation; no implemented
bootstrap or significance analysis is claimed here.

\paragraph{Model, caps, and concurrency.}
All three optimization roles use Azure Foundry \texttt{DeepSeek-V4-Flash},
temperature 0, top-$p$ 0.95, and an output-token cap of 8,192. The provider's
default reasoning mode is enabled, without guaranteeing returned reasoning.
Per-run limits are 30 iterations, depth three, 3,600 seconds, and a \$2 usage
budget. Each validation subject, including the incumbent, has a \$100 cumulative
budget. There are three mining run workers. Validation allows five subject
workers but has only two measured subjects, with one run worker per subject;
recursive calls can add provider concurrency. Final evaluation has a separate
run-worker option and uses three attempts per task per condition. Record the
chosen worker count in the evaluation manifest; it does not change repetition
counts.

The main comparison explicitly selects \texttt{initial} and \texttt{sh\_rlm}.
The default reference grid is \texttt{b1}, \texttt{h0\_star},
\texttt{lambda\_rlm}, \texttt{sh\_rlm}. The latter's typed baseline is identified
by its method kind, reconstruction version, upstream revision, and source hash;
it is not an unmodified released OOLONG-Pairs implementation. Model/deployment,
rendered prompts, source/configuration snapshots, split IDs/hashes, exact edit
payloads, harness lineage, and analysis versions belong in the study artifacts.
Freezing these before a run does not by itself constitute preregistration.
Previously inspected development/test material must retain its exposure history
if later used to support a result.

\section{Measurement, Cost, and Saved Observations}
\label{app:measurement}
\label{sec:cost}

\paragraph{Quality and surface activity.}
Exact pass is set equality, not rounded F1 equal to one. OOLONG-Pairs stores
per-attempt precision, recall, and F1 to three decimals; reported mean F1 averages
recorded attempt values, including verifier-declared zeros for completed runtime,
resource, format, and content-filter failures. It is macro averaging across
attempts, not pooled pair-level micro-F1. Missing/extra counts are reported only
where available, with coverage, rather than filling absent diagnostics with zero.
Other environments supply their own metric definitions or remain unassessed.

Surface activity distinguishes submitted, admitted, evaluated, and promoted
edits. Constituents inherit promotion membership from a promoted batch without
individual efficacy estimates. Per-round quality plots use the freshly measured
candidate on promotion and freshly measured incumbent otherwise; a round without
a measurement is missing, not a zero or a carried-forward score. Environment
labels come from the configuration. Trace summaries can report depth, call
counts, and the implemented whole-input-subcall proxy: maximum descendant prompt
characters divided by root prompt characters exceeds 0.8. Its denominator is
walkable runs, passing and failing, with unwalkable runs separate; it is neither
a token ratio nor proof of meaningless decomposition.

\paragraph{Cost.}
An ordinary completed round with an evaluable batch schedules
\begin{equation}
N_{\mathrm{round}}=m n_{\mathrm{in}}+2v n_{\mathrm{ho}}.
\label{eq:round-cost}
\end{equation}
The current profile gives $40+20=60$ task attempts, independent of how many
edits the batch contains, or 180 over three such rounds. No evaluable batch
means no new validation leg. With 10 short and 40 long tasks, final evaluation
requires $3(10+40)J=150J$ attempts for $J$ conditions: 300 for initial versus
final, 600 for the default four-condition grid, or 750 for all five. These are
scheduling counts, not measured results or provider-call counts. Root attempts
may issue many recursive/model requests. Attribution, proposal, repairs, and
operational retries also consume resources.

Configured accounting rates are \$0.19 per million input tokens and \$0.51 per
million output tokens, as recorded for this study; they are not a current market
quote. Actual usage and its coverage will be reported separately from scheduling
projections. Task usage is not added again to a stage total
that already includes it. Breaker reservations are not billed spend; interrupted
or unreported requests can leave a lower bound, and the per-run budget is not a
guarantee of complete invoice accounting.

\paragraph{Provider-returned observations.}
Instrumented execution, mining, validation, final evaluation, attribution,
proposal/repair, recursive calls, compaction, and fallback paths save returned
response/reasoning observations separately from executable code and final
answers. References carry call IDs, paths, and hashes. Run traces point to
\path{<run_id>.llm_calls/<call_id>/}; attribution and proposal artifacts point to
their stage-specific \path{llm_calls/} directories. Started, response, and finished
events distinguish failed or interrupted calls. Cache hits retain source
provenance and do not add spend.

Availability is explicit: \texttt{returned}, \texttt{opaque\_only},
\texttt{not\_returned}, \texttt{unsupported\_client},
\texttt{legacy\_unavailable}, or \texttt{no\_response}. Returned content can be
text or a summary; this is not guaranteed access to full private reasoning.
Unknown reasoning-token counts stay unknown and are not added to reported output
usage again. Legacy traces cannot be backfilled with unreturned reasoning.
Observation artifacts are not automatically fed into attribution digests,
proposal evidence, or history. Ordinary assistant prose in a digest is a
different source. Observation persistence errors remain visible failures.

% AAAI: full-width table* (the version strings are too long for one column).
\begin{table*}[t]
\centering\small
\begin{tabularx}{\textwidth}{@{}Xl@{}}
\toprule
Contract & Version \\
\midrule
Taxonomy / capabilities & 3.3.0 / \texttt{surface-capabilities/v1} \\
Attribution prompt / validator / digest & 1.6.0 / 1.2.0 / 1.8.0 \\
Proposal prompt / validator & 4.3.0 / 4.2.0 \\
Proposal response / text & \texttt{proposal-selection/v2} / \texttt{literal-text/v1} \\
Evidence selector / diagnostic history & 4.3.0 / 2.1.0 \\
Compact history & \texttt{proposal-history/v3} \\
Validation / summary & \texttt{heldout-batch/v2} / \texttt{shrlm-validation-summary/v3} \\
Behavior / provider observations & \texttt{behavior-observations/v1} / \texttt{rlm-llm-observation/v1} \\
\bottomrule
\end{tabularx}
\caption{Contracts at the documented revision; archive the versions actually
used by the forthcoming run with its results.}
\end{table*}

\section{Optional Follow-up Studies}
\label{app:extensions}
\label{app:finetune}
\label{app:alignment}

Cross-environment transfer would require a separately specified target protocol
and completed evaluation. A weight-training comparison would require an available
trainable checkpoint, justified backbone matching, and measured compute; no
fine-tuning condition or hardware allocation is claimed for this hosted study.
Joint batch outcomes suffice for the present design. An optional future
leave-one-surface-out study would revert a final surface to its starting value
and rerun the same evaluation; it would measure that final-state counterfactual,
not the marginal contribution of every historical edit overwritten along the way.
No such ablation has been completed here.

An optional safety probe could compare frozen harness states for changes in
refusal and harmful-answer behavior, motivated by work on
misevolution~\citep{shao2025misevolution}. The task gate does not measure these
outcomes, so improved benchmark behavior would not establish safe deployment.
Any probe would need its own data, scoring, and budget. Textual reasoning
instructions are editable; provider reasoning settings are experiment owned.

% Reproducibility checklist: if ReproducibilityChecklist.tex has its own
% \documentclass, compile and submit it separately instead of \input-ing it.
% \makeatletter
\@ifundefined{isChecklistMainFile}{
  % We are compiling a standalone document
  \newif\ifreproStandalone
  \reproStandalonetrue
}{
  % We are being \input into the main paper
  \newif\ifreproStandalone
  \reproStandalonefalse
}
\makeatother

\ifreproStandalone
\documentclass[letterpaper]{article}
\usepackage[submission]{aaai2027}
\setlength{\pdfpagewidth}{8.5in}
\setlength{\pdfpageheight}{11in}
\frenchspacing

\begin{document}
\fi
\setlength{\leftmargini}{20pt}
\makeatletter\def\@listi{\leftmargin\leftmargini \topsep .5em \parsep .5em \itemsep .5em}
\def\@listii{\leftmargin\leftmarginii \labelwidth\leftmarginii \advance\labelwidth-\labelsep \topsep .4em \parsep .4em \itemsep .4em}
\def\@listiii{\leftmargin\leftmarginiii \labelwidth\leftmarginiii \advance\labelwidth-\labelsep \topsep .4em \parsep .4em \itemsep .4em}\makeatother

\setcounter{secnumdepth}{0}
\renewcommand\thesubsection{\arabic{subsection}}
\renewcommand\labelenumi{\thesubsection.\arabic{enumi}}

\newcounter{checksubsection}
\newcounter{checkitem}[checksubsection]

\newcommand{\checksubsection}[1]{%
  \refstepcounter{checksubsection}%
  \paragraph{\arabic{checksubsection}. #1}%
  \setcounter{checkitem}{0}%
}

\newcommand{\checkitem}{%
  \refstepcounter{checkitem}%
  \item[\arabic{checksubsection}.\arabic{checkitem}.]%
}
\newcommand{\question}[2]{\normalcolor\checkitem #1 #2 \color{blue}}
\newcommand{\ifyespoints}[1]{\makebox[0pt][l]{\hspace{-15pt}\normalcolor #1}}

\section*{Reproducibility Checklist}

\vspace{1em}
\hrule
\vspace{1em}

\textbf{Instructions for Authors:}

This document outlines key aspects for assessing reproducibility. Please provide your input by editing this \texttt{.tex} file directly.

For each question (that applies), replace the ``Type your response here'' text with your answer.

\vspace{1em}
\noindent
\textbf{Example:} If a question appears as
%
\begin{center}
\noindent
\begin{minipage}{.9\linewidth}
\ttfamily\raggedright
\string\question \{Proofs of all novel claims are included\} \{(yes/partial/no)\} \\
Type your response here
\end{minipage}
\end{center}
you would change it to:
\begin{center}
\noindent
\begin{minipage}{.9\linewidth}
\ttfamily\raggedright
\string\question \{Proofs of all novel claims are included\} \{(yes/partial/no)\} \\
yes
\end{minipage}
\end{center}
%
Please make sure to:
\begin{itemize}\setlength{\itemsep}{.1em}
\item Replace ONLY the ``Type your response here'' text and nothing else.
\item Use one of the options listed for that question (e.g., \textbf{yes}, \textbf{no}, \textbf{partial}, or \textbf{NA}).
\item \textbf{Not} modify any other part of the \texttt{\string\question} command or any other lines in this document.\\
\end{itemize}

You can \texttt{\string\input} this .tex file right before \texttt{\string\end\{document\}} of your main file or compile it as a stand-alone document. Check the instructions on your conference's website to see if you will be asked to provide this checklist with your paper or separately.

\vspace{1em}
\hrule
\vspace{1em}

% The questions start here

\checksubsection{General Paper Structure}
\begin{itemize}

\question{Includes a conceptual outline and/or pseudocode description of AI methods introduced}{(yes/partial/no/NA)}
Type your response here

\question{Clearly delineates statements that are opinions, hypothesis, and speculation from objective facts and results}{(yes/no)}
Type your response here

\question{Provides well-marked pedagogical references for less-familiar readers to gain background necessary to replicate the paper}{(yes/no)}
Type your response here

\end{itemize}
\checksubsection{Theoretical Contributions}
\begin{itemize}

\question{Does this paper make theoretical contributions?}{(yes/no)}
Type your response here

	\ifyespoints{\vspace{1.2em}If yes, please address the following points:}
        \begin{itemize}
	
	\question{All assumptions and restrictions are stated clearly and formally}{(yes/partial/no)}
	Type your response here

	\question{All novel claims are stated formally (e.g., in theorem statements)}{(yes/partial/no)}
	Type your response here

	\question{Proofs of all novel claims are included}{(yes/partial/no)}
	Type your response here

	\question{Proof sketches or intuitions are given for complex and/or novel results}{(yes/partial/no)}
	Type your response here

	\question{Appropriate citations to theoretical tools used are given}{(yes/partial/no)}
	Type your response here

	\question{All theoretical claims are demonstrated empirically to hold}{(yes/partial/no/NA)}
	Type your response here

	\question{All experimental code used to eliminate or disprove claims is included}{(yes/no/NA)}
	Type your response here
	
	\end{itemize}
\end{itemize}

\checksubsection{Dataset Usage}
\begin{itemize}

\question{Does this paper rely on one or more datasets?}{(yes/no)}
Type your response here

\ifyespoints{If yes, please address the following points:}
\begin{itemize}

	\question{A motivation is given for why the experiments are conducted on the selected datasets}{(yes/partial/no/NA)}
	Type your response here

	\question{All novel datasets introduced in this paper are included in a data appendix}{(yes/partial/no/NA)}
	Type your response here

	\question{All novel datasets introduced in this paper will be made publicly available upon publication of the paper with a license that allows free usage for research purposes}{(yes/partial/no/NA)}
	Type your response here

	\question{All datasets drawn from the existing literature (potentially including authors' own previously published work) are accompanied by appropriate citations}{(yes/no/NA)}
	Type your response here

	\question{All datasets drawn from the existing literature (potentially including authors' own previously published work) are publicly available}{(yes/partial/no/NA)}
	Type your response here

	\question{All datasets that are not publicly available are described in detail, with explanation why publicly available alternatives are not scientifically satisficing}{(yes/partial/no/NA)}
	Type your response here

\end{itemize}
\end{itemize}

\checksubsection{Computational Experiments}
\begin{itemize}

\question{Does this paper include computational experiments?}{(yes/no)}
Type your response here

\ifyespoints{If yes, please address the following points:}
\begin{itemize}

	\question{This paper states the number and range of values tried per (hyper-) parameter during development of the paper, along with the criterion used for selecting the final parameter setting}{(yes/partial/no/NA)}
	Type your response here

	\question{Any code required for pre-processing data is included in the appendix}{(yes/partial/no)}
	Type your response here

	\question{All source code required for conducting and analyzing the experiments is included in a code appendix}{(yes/partial/no)}
	Type your response here

	\question{All source code required for conducting and analyzing the experiments will be made publicly available upon publication of the paper with a license that allows free usage for research purposes}{(yes/partial/no)}
	Type your response here
        
	\question{All source code implementing new methods have comments detailing the implementation, with references to the paper where each step comes from}{(yes/partial/no)}
	Type your response here

	\question{If an algorithm depends on randomness, then the method used for setting seeds is described in a way sufficient to allow replication of results}{(yes/partial/no/NA)}
	Type your response here

	\question{This paper specifies the computing infrastructure used for running experiments (hardware and software), including GPU/CPU models; amount of memory; operating system; names and versions of relevant software libraries and frameworks}{(yes/partial/no)}
	Type your response here

	\question{This paper formally describes evaluation metrics used and explains the motivation for choosing these metrics}{(yes/partial/no)}
	Type your response here

	\question{This paper states the number of algorithm runs used to compute each reported result}{(yes/no)}
	Type your response here

	\question{Analysis of experiments goes beyond single-dimensional summaries of performance (e.g., average; median) to include measures of variation, confidence, or other distributional information}{(yes/no)}
	Type your response here

	\question{The significance of any improvement or decrease in performance is judged using appropriate statistical tests (e.g., Wilcoxon signed-rank)}{(yes/partial/no)}
	Type your response here

	\question{This paper lists all final (hyper-)parameters used for each model/algorithm in the paper’s experiments}{(yes/partial/no/NA)}
	Type your response here


\end{itemize}
\end{itemize}
\ifreproStandalone
\end{document}
\fi

\end{document}